\section{Mathematical preliminaries and proof convention}
\label{sec:preliminaries}

\subsection{Functions, images, and fibers}

The notation for functions, images, fibers, composition, inverses, and equivalence relations follows standard set-theoretic usage \parencite[Chs. 8--11]{Halmos1974}. For a function \(f:X\to Y\), the image of a subset \(A\subseteq X\) is
\[
f[A]=\{f(x):x\in A\}.
\]
The range is \(\im(f)=f[X]\). The fiber over \(y\in Y\) is
\[
f^{-1}(y)=\{x\in X:f(x)=y\}.
\]
The function is injective when \(f(x)=f(x')\) implies \(x=x'\). It is surjective when every \(y\in Y\) has at least one preimage.

A function \(g:Y\to X\) is a left inverse of \(f\) when
\[
g\circ f=\id_X.
\]
A left inverse can exist only when \(f\) is injective \parencite[Chs. 1 and 4]{MacLaneBirkhoff1999}. This elementary fact is the exact-decoding obstruction used throughout the paper.

\subsection{Equivalence relations and quotients}

An equivalence relation partitions a set into disjoint equivalence classes. The quotient \(X/{\sim}\) is the set of those classes. The canonical class map is
\[
q:X\to X/{\sim},
\qquad
q(x)=[x].
\]
A function on \(X\) descends to the quotient precisely when it is constant on each equivalence class. The descended function is then unique because every quotient class has a representative. This is the elementary set-level form of a quotient universal property \parencite{MacLaneBirkhoff1999,Leinster2014}.

\subsection{State-level and model-level notation}

The symbol \(X\) denotes a source or world state space. The symbol \(O\) denotes an accessible observation space. The observation map is \(\pi:X\to O\). A generative model is written
\[
\Model=(X_\Model,\pi_\Model).
\]
Its reachable observation set is \(\Obs(\Model)\). The possible outcome set of an admissible experiment \(e\) is \(\Out_e(\Model)\).

Two equivalence relations are used. State-level observational equivalence identifies \(x,y\in X\) when \(\pi(x)=\pi(y)\). Model-level empirical equivalence identifies \(\Model\) and \(\Model'\) when \(\Obs(\Model)=\Obs(\Model')\). The levels are not interchangeable.

\subsection{Logical status of statements}

Definitions fix the objects and relations used in the analysis. Propositions, lemmas, corollaries, and theorems are established by the proofs that follow them. Finite computations evaluate explicitly enumerated instances of those definitions. Interpretations, hypotheses, applications, and future research questions are stated separately and do not inherit proof from the mathematical results. Literature citations identify established precedent, comparison, or empirical context.
